\section{Validação do protótipo}
\label{sec:validacao}

\subsection{Conformidade com a especificação}

A validação tem três níveis. Os \textbf{testes automatizados}
(Secção~\ref{sec:prototipo}) verificam cada regra de aceitação, os vetores de
teste do envelope (partilhados pela implementação em C++ e pela em Python, que
assim não podem divergir) e os vetores do RFC~8439 para o ChaCha20-Poly1305. Os
\textbf{cenários} (Secção~\ref{sec:protecao}) verificam o sistema completo contra
um \emph{broker} real: os 39 cenários deram o resultado previsto no
Raspberry~Pi, na própria campanha (\texttt{results/pi5-20261113/scenarios.md}). No
mesmo \emph{commit} da campanha (\texttt{79c4878}), os 93 testes em C++ e os 61
testes em Python do laboratório passaram no Raspberry~Pi
(\texttt{results/validation/pi5-20261113/}); só os sete testes que dependem do
\texttt{tshark} não correram, por o Wireshark não estar instalado no
Raspberry~Pi.
A \textbf{matriz de ataques} confirma que cada configuração impede o que o desenho
prevê (Tabela~\ref{tab:protecao}), com as três repetições de acordo em todas as
células.

\subsection{Validação do instrumento de medição}

Os pilotos no Raspberry~Pi validaram também o próprio laboratório, e revelaram
dois defeitos que teriam falseado a campanha:

\begin{itemize}[leftmargin=*]
  \item \textbf{Deriva dos relógios.} No primeiro piloto, a latência mediana de K0
    saiu negativa ($-455~\mu$s) e variava de 300 a 1\,055~$\mu$s entre repetições
    da mesma configuração. O desvio entre os relógios crescia de forma regular
    (de 32,9 para 51,8~ms ao longo da subida de carga, cerca de 13~ppm), e a
    estimação usava a melhor amostra do último minuto, que podia estar 0,8~ms
    atrasada. Com a reta ajustada às amostras
    (Secção~\ref{sec:metodo}), a latência mediana de K0 ficou entre 132 e
    133~$\mu$s nas três repetições do terceiro piloto. Um teste reproduz a
    deriva observada e exige um erro inferior a 50~$\mu$s.
  \item \textbf{Contaminação entre patamares.} No segundo piloto, uma execução
    de K8 no primeiro patamar (500~mensagens/s) registou 77\,388 rejeições e p99
    de 1,3~s: seguia-se a um patamar de K11 em sobrecarga, cujas mensagens retidas
    chegaram depois da mudança de chave e de limite de débito. Se a perda tivesse
    passado de 1\,\%, a subida de carga de K8 teria parado logo nesse patamar e o
    seu débito máximo ficaria por medir. Com a espera
    pelo escoamento, nenhuma execução posterior a uma saturação teve rejeições no
    terceiro piloto.
\end{itemize}

\subsection{Calibração da validade temporal}

A Parte~1 fixou valores iniciais, $\Delta_a = 60$~s e $\Delta_s = 5$~s, a calibrar
nesta parte a partir da latência medida e do desvio de relógio observado. Cada
execução regista a cauda da latência e o desvio entre o relógio de cada máquina de
sensores e o do \emph{gateway}.

Na campanha, a maior latência de uma mensagem legítima foi de 43,2~ms (p99,9
máximo de 37,1~ms), contando as execuções de custo e os patamares da subida de
carga que cada configuração sustenta. Com o NTP desligado, o desvio entre os
relógios das duas máquinas cresceu até 289~ms ao longo das cinco horas; depois da
correção pelos agentes, a incerteza residual nunca passou de 0,21~ms. Com margem
de segurança~2 e arredondamento a 100~ms, os valores propostos são:

\begin{itemize}[leftmargin=*]
  \item \textbf{relógios sincronizados} (NTP ou PTP nos dispositivos):
    $\Delta_s = 100$~ms e $\Delta_a = 100$~ms;
  \item \textbf{relógios sem sincronização}, em que $\Delta_s$ tem de cobrir o
    desvio observado: $\Delta_s = 600$~ms e $\Delta_a = 700$~ms.
\end{itemize}

Os dois casos estão muito abaixo dos 60~s e 5~s iniciais da Parte~1: uma mensagem
antiga só é aceite se tiver menos de um segundo, em vez de um minuto. O segundo
caso depende de quanto tempo os relógios ficam sem acerto (aqui, cinco horas a
cerca de 16~ppm); num dispositivo sem relógio de tempo real que reinicia, o desvio
pode ser arbitrário, e é a sincronização que o resolve, não um $\Delta_s$ maior.
Para as ordens a atuadores, em que uma ordem atrasada é ela própria um risco, a
validade curta é a que interessa: com relógios sincronizados, uma ordem com mais
de 100~ms é recusada. Por fluxo (\texttt{FreshnessRule}), a telemetria pode
manter uma validade mais larga sem afetar as ordens.

\subsection{Reinício do recetor}

O recetor guarda em disco a última época de chaves que aceitou. Depois de um
reinício, a janela anti-repetição está vazia, mas só são aceites chaves de uma
época posterior: as mensagens da época anterior, incluindo repetições capturadas
antes do reinício, são rejeitadas como repetição, e o tráfego retoma quando os
emissores passam à época seguinte. Um teste de aceitação verifica esta sequência
(\texttt{AfterARestartOnlyALaterEpochIsAccepted}), e outro que uma falha a ler ou
a escrever a época recusa as mensagens em vez de as aceitar. No laboratório, o
\emph{gateway} corre sempre com esta época em disco.
O tempo sem serviço depois de um reinício é, por isso, o tempo que a implantação
leva a passar os emissores à época seguinte. Depende da distribuição de chaves,
fora do âmbito deste trabalho (Secção~\ref{sec:ameacas}), e não foi medido.

\subsection{Rotação e revogação}

Os dois cenários passaram no Raspberry~Pi. Em \texttt{key-rotation}, todas as
chaves do conteúdo são substituídas com o tráfego a correr, sem nenhuma rejeição e
com perda de 0,1\,\%. Em \texttt{revocation}, o sensor~1 é revogado (chaves e
palavra-passe): fica de fora e os outros 49 continuam, com o débito entregue a
descer de 1\,000 para 980~mensagens/s, exatamente a parte do sensor revogado.
